% Unit 18 teaching — Sneaky Symmetry \& Sneaky Product Rule (Yang Ch.2)
\input{technique-note.tex}
\title{Integration Bee --- Unit 18:\\ Sneaky Symmetry \& Sneaky Product Rule}
\author{UK Uni Integration Bee prep\\ Yang Topics Guide 2nd Ed §§2.9 Sneaky Symmetry, 2.10 Sneaky Product Rule}
\date{}
\begin{document}
\maketitle

\section*{Status}
Yang Chapter 2 concept unit (easy drills only; no past papers).
Tags: \texttt{yang:2.9}, \texttt{yang:2.10}, \texttt{sneaky-symmetry}, \texttt{sneaky-product-rule}.

\section{Official anchors}
\begin{itemize}
  \item Yang Topics Guide 2nd Ed §§2.9 Sneaky Symmetry, 2.10 Sneaky Product Rule
        (\texttt{sources/user/Integration\_Bee\_Topics\_Guide\_Yang\_2nd\_Ed.pdf}).
  \item Syllabus §1 techniques support
        (\texttt{sources/UK\_University\_Integration\_Bee\_Syllabus.pdf}).
\end{itemize}

\section{Learning goals}
\begin{enumerate}
  \item Center an ugly product with $u=x-a$ so an odd integrand cancels.
  \item Spot $f'g+fg'$ hiding as a sum of scary-looking terms.
  \item Differentiate proposed answers mentally before grinding parts.
\end{enumerate}

\section{Pedagogy}
Worked $\to$ faded $\to$ interleaved easy drills. Chat Flow 1.

\section{Worked example --- \texttt{U18-D01} (\texttt{yang:2.9})}
\[
\int_{2}^{4}(x-1)(x-2)(x-3)(x-4)(x-5)\,dx.
\]
$u=x-3$ makes an odd integrand on $[-1,1]$ $\Rightarrow 0$.

\section{Worked example --- \texttt{U18-D05} (\texttt{yang:2.10})}
\[
\int\bigl(\sec^{3}x+\sec x\tan^{2}x\bigr)\,dx.
\]
Rewrite as $\sec\cdot\sec^{2}+(\sec\tan)\tan=\frac{d}{dx}(\sec\tan)$.

\section{Faded example}
Try U18-D07: $\int e^{x}(\sin x-\cos x)\,dx=-\,e^{x}\cos x+C$.

\section{Drill pointers}
See \texttt{drills.tex} / \texttt{drills.md}. SymPy feed: \texttt{verify/drills\_unit18.json}.

\end{document}
